% Part: lambda-calculus
% Chapter: introduction
% Section: computable-lambda

\documentclass[../../../include/open-logic-section]{subfiles}

\begin{document}

\olfileid{lam}{int}{lrp}
\olsection{গণনসাধ্য অপেক্ষকগুলি \usetoken{S}{lambda definable}}

\begin{thm}
\ollabel{thm:computable-lambda}
প্রতিটি আংশিক গণনসাধ্য অপেক্ষক !!{lambda definable}।
\end{thm}

\begin{proof}
দেখাতে হবে যে প্রতিটি আংশিক গণনসাধ্য অপেক্ষক~$f$-কে কোনো ল্যাম্বডা
পদ~$F$ !!{lambda defined} করে। ক্লিনির স্বাভাবিক রূপ উপপাদ্য অনুসারে,
প্রথমে প্রতিটি আদিম পুনরাবৃত্ত অপেক্ষককে কোনো ল্যাম্বডা পদ
!!{lambda defined} করে এবং তারপর !!{lambda definable} অপেক্ষকগুলি
উপযুক্ত মিশ্রণ ও সীমাহীন অনুসন্ধানের অধীনে বদ্ধ---এটুকু দেখালেই যথেষ্ট। প্রতিটি আদিম
পুনরাবৃত্ত অপেক্ষককে কোনো ল্যাম্বডা পদ !!{lambda defined} করে দেখাতে,
প্রারম্ভিক অপেক্ষকগুলি !!{lambda definable} এবং !!{lambda definable}
আংশিক অপেক্ষকগুলি মিশ্রণ, আদিম পুনরাবৃত্তি ও সীমাহীন অনুসন্ধানের অধীনে
বদ্ধ---এটুকু দেখালেই যথেষ্ট।
\end{proof}

প্রমাণের বাকি অংশ আরও সহজে পড়ার জন্য আমরা একটি বেশি প্রচলিত সংকেতলিপি
ব্যবহার করব। যেমন, $Mxyz$-এর বদলে $M(x, y, z)$ লিখব। এই লিপি অর্থের
ইঙ্গিত দিলেও মনে রাখতে হবে, টাইপবিহীন ল্যাম্বডা ক্যালকুলাসের পদগুলির সঙ্গে
কোনো নির্দিষ্ট স্থানসংখ্যা যুক্ত নেই; তাই একই পদ~$M$-এর জন্য $M(x,y)$
লেখা যতটা অর্থপূর্ণ, $M(x,y,z,w)$ লেখাও ততটাই অর্থপূর্ণ। তবে এই লিপি
ব্যবহার করলে বোঝা যায় যে আমরা~$M$-কে তিন চলরাশির অপেক্ষক হিসেবে ধরছি;
এবং সংজ্ঞাগুলির উদ্দেশ্যও আরও স্পষ্ট হয়। একইভাবে, ``$M =
\lambd[x][\lambd[y][\lambd[z][\dots]]]$ দ্বারা $M$-কে সংজ্ঞায়িত করো''
বলার বদলে বলব, ``$M(x,y,z) = \dots$ দ্বারা $M$-কে সংজ্ঞায়িত করো''।

\end{document}
